\section{Background: definitions, surveys and the boundary question}\label{sec:background}

When does an attitude about Israel express hostility toward Jews, and when is it a political attitude like any other? Part of this \emph{boundary question} is normative: what ought to count as antisemitism. Part is empirical: how attitudes about Jews and attitudes about Israel are in fact structured in a population. This section reviews the three definitions that dominate current debate (\cref{sec:bg-definitions}), the main survey programmes that measure antisemitic attitudes (\cref{sec:bg-surveys}), and prior empirical work on the relationship between anti-Israel and antisemitic attitudes (\cref{sec:bg-prior}). It then explains why measurement invariance (\cref{sec:bg-invariance}) and indirect questioning (\cref{sec:bg-list}) are needed to answer the empirical part.

We follow two conventions. First, we do not adjudicate between definitions. We describe each in its own terms, and we identify the empirical premises that the definitions appear to rest on, because those premises are what our methods can test. Second, every survey figure quoted below is attributed to the organisation or authors that published it. None of them is an input to the simulation study of \cref{sec:simulation}, and we have fielded no survey.

\subsection{Three definitions and the claims embedded in them}\label{sec:bg-definitions}

\paragraph{The IHRA working definition (2016)}
On 26 May 2016 the Plenary of the International Holocaust Remembrance Alliance (IHRA), meeting in Bucharest, adopted a ``non-legally binding working definition'' of antisemitism \citep{ihra2016}. Its core reads: ``Antisemitism is a certain perception of Jews, which may be expressed as hatred toward Jews.'' A second sentence states that rhetorical and physical manifestations of antisemitism are directed toward Jewish or non-Jewish individuals or their property, and toward Jewish community institutions and religious facilities. The text then states that the examples which follow ``may serve as illustrations'' to guide IHRA in its work. It says that manifestations might include the targeting of the State of Israel, conceived as a Jewish collectivity. It adds: ``However, criticism of Israel similar to that leveled against any other country cannot be regarded as antisemitic.'' Eleven contemporary examples follow. They are introduced as instances that could, ``taking into account the overall context'', count as antisemitism. Several concern Israel. They include denying the Jewish people their right to self-determination (for example, by claiming that the existence of a State of Israel is a racist endeavour); applying double standards by requiring of Israel a behaviour not expected or demanded of any other democratic nation; using symbols and images associated with classic antisemitism to characterise Israel or Israelis; drawing comparisons of contemporary Israeli policy to that of the Nazis; and holding Jews collectively responsible for actions of the State of Israel \citep{ihra2016}.

\paragraph{The Jerusalem Declaration on Antisemitism (2021)}
The Jerusalem Declaration on Antisemitism (JDA) was published in March 2021 by scholars working, in its own words, in antisemitism studies and related fields, including Jewish, Holocaust, Israel, Palestine and Middle East studies. It had 210 initial signatories \citep{jda2021}. It defines antisemitism as ``discrimination, prejudice, hostility or violence against Jews as Jews (or Jewish institutions as Jewish)''. Fifteen guidelines follow in three groups: general guidelines (A); examples relating to Israel and Palestine that, ``on the face of it'', are antisemitic (B); and examples that, on the face of it, are not antisemitic, whether or not one approves of the view or action (C). Group~B includes applying the symbols, images and negative stereotypes of classical antisemitism to the State of Israel, holding Jews collectively responsible for Israel's conduct, and denying the right of Jews in Israel to exist and flourish, collectively and individually, as Jews. Group~C includes supporting Palestinian rights, criticising or opposing Zionism as a form of nationalism, evidence-based criticism of Israel as a state, and boycott, divestment and sanctions as forms of political protest. It also states that it is not antisemitic, ``in and of itself'', to compare Israel with other historical cases, including settler-colonialism or apartheid. The JDA says that context, including the intention behind an utterance, a pattern of speech over time and the identity of the speaker, can matter. Its authors describe the IHRA definition as unclear in key respects and open to different interpretations, and they present the JDA as an alternative to it or as a corrective for bodies that have adopted it \citep{jda2021}.

\paragraph{The Nexus Document (2021)}
The Nexus Document was produced by the Nexus Task Force, initially hosted by the Knight Program on Media and Religion at the University of Southern California's Annenberg School. It was published in February 2021 and updated in June 2024 \citep{nexus2021}. It states that ``antisemitism consists of anti-Jewish beliefs, attitudes, actions or systemic conditions''. It then sets out when Israel-related expression is and is not antisemitic. Expression is antisemitic, for example, when it promotes myths, stereotypes or attitudes about Zionism or Israel that derive from or reinforce antisemitic accusations and tropes. As a general rule, criticism of Zionism and Israel, opposition to Israel's policies, and nonviolent political action directed at Israel or its policies ``should not, as such, be deemed antisemitic''. Treating Israel differently from other countries because it is a Jewish state is listed as antisemitic, but paying disproportionate attention to Israel, or treating Israel differently from other countries, is not in itself treated as proof of antisemitism \citep{nexus2021}.

\begin{table}[tbp]
\centering
\caption{Three definitions of antisemitism, summarised from their own texts. ``Israel-related content'' lists only provisions quoted or paraphrased in \cref{sec:bg-definitions}; it is not a complete inventory.}\label{tab:bg-definitions}
\footnotesize
\setlength{\tabcolsep}{4pt}
\renewcommand{\arraystretch}{1.15}
\begin{tabularx}{\textwidth}{@{}>{\raggedright\arraybackslash}p{2.0cm}>{\raggedright\arraybackslash}p{2.6cm}>{\raggedright\arraybackslash}X>{\raggedright\arraybackslash}X@{}}
\toprule
Document & Issuer; date & Core definition & Israel-related content \\
\midrule
IHRA working definition & IHRA Plenary, Bucharest; 26 May 2016; non-legally binding & ``A certain perception of Jews, which may be expressed as hatred toward Jews'' & Eleven illustrative examples, several concerning Israel (self-determination, double standards, classic imagery, Nazi comparisons, collective responsibility), applicable ``taking into account the overall context''; criticism similar to that of any other country is not antisemitic \\
\addlinespace
Jerusalem Declaration (JDA) & Scholars; March 2021; 210 initial signatories & ``Discrimination, prejudice, hostility or violence against Jews as Jews (or Jewish institutions as Jewish)'' & Guidelines 6--10: examples antisemitic on the face of it (e.g.\ classic stereotypes applied to Israel, collective responsibility); guidelines 11--15: examples not antisemitic on the face of it (e.g.\ opposing Zionism as nationalism, BDS, comparisons with settler-colonialism or apartheid) \\
\addlinespace
Nexus Document & Nexus Task Force; February 2021, updated June 2024 & ``Anti-Jewish beliefs, attitudes, actions or systemic conditions'' & Antisemitic when drawing on antisemitic myths, stereotypes or tropes, or when treating Israel differently because it is a Jewish state; criticism, opposition to policies and nonviolent political action ``should not, as such, be deemed antisemitic''; disproportionate attention not in itself proof \\
\bottomrule
\end{tabularx}
\end{table}

\paragraph{Agreement and disagreement}
\Cref{tab:bg-definitions} summarises the three texts. They share some ground. All three list the transfer of classic antisemitic imagery or tropes to Israel among their examples of antisemitism, and all three list, in similar terms, holding Jews collectively responsible for Israel's conduct. All three state, in different words and with different qualifications, that criticism of Israel is not in itself antisemitic; the IHRA text refers to criticism ``similar to that leveled against any other country''. They differ on several categories of Israel-related speech. These include claims that deny the legitimacy of a Jewish state or oppose Zionism; double standards and disproportionate attention; comparisons of Israel with Nazi Germany, with apartheid or with settler-colonialism; and boycott campaigns. The IHRA text lists some of these among its contextual examples. The JDA lists some of them among its examples that are not antisemitic on the face of it. The Nexus Document treats several as not antisemitic as such, but antisemitic when they draw on antisemitic tropes. They also differ in how much weight they give to context and intent \citep{ihra2016,jda2021,nexus2021}.

\paragraph{Normative and empirical components}
Much of this disagreement is normative: it concerns which speech acts ought to be called antisemitic, by whom, and with what consequences. Survey data cannot settle it. But each position also rests partly on claims about how attitudes are distributed in a population. One such claim is that endorsement of a given Israel-related statement typically reflects the same hostility toward Jews that drives endorsement of classic stereotypes. The competing claim is that, in the population at large, endorsement of that statement is mainly driven by a distinct political attitude. These are claims about covariance structure, and they can be stated as restrictions on a measurement model. \Cref{def:boundary} gives the item-level estimand that separates them, and \cref{sec:model,sec:identification} give the conditions under which it is identified. The estimand describes response patterns in a population; it does not classify any utterance, speaker or context. A high general-factor share for an item does not make each endorsement of it antisemitic, and a low share does not make the item's content acceptable in every context. We return to this distinction in \cref{sec:ethics}.

\subsection{Survey measurement of antisemitic attitudes}\label{sec:bg-surveys}

\paragraph{Stereotype-count indices (ADL)}
The Anti-Defamation League (ADL) measures antisemitic attitudes with a fixed battery of eleven statements expressing negative stereotypes about Jews. In its Global 100 survey (2014), respondents who answered ``probably true'' to at least six of the eleven were classified as harbouring antisemitic attitudes \citep{adl2014global}. ADL's United States surveys use the same eleven tropes, which it reports having measured since 1964, and the same threshold of six or more \citep{adl2023attitudes}. Its 2023 report on the United States also included statements about Israel. ADL describes belief in anti-Jewish tropes and anti-Israel sentiment as distinct but overlapping in substantial ways \citep{adl2023attitudes}. Indices of this kind are transparent and comparable over time when the battery is held fixed. Their group comparisons, however, rest on an assumption that each statement means the same thing, and has the same threshold for agreement, in every group compared. \Cref{prop:index} shows what happens to a count index when that assumption fails. We make no claim about the results of any real index.

\paragraph{The elastic view (JPR)}
The Institute for Jewish Policy Research (JPR), in partnership with the Community Security Trust, published a study of attitudes towards Jews and Israel in Great Britain in 2017 \citep{staetsky2017}. It drew on a face-to-face sample and an online-panel sample. The report proposes an ``elastic view'' of antisemitism: as an attitude, antisemitism exists at different levels of intensity, so no single figure captures its level in a society. The report distinguishes \emph{counting antisemites}, which requires a threshold, from \emph{measuring antisemitism}, which tracks how widely antisemitic ideas are diffused. It measures anti-Israel attitudes separately from antisemitic attitudes. It reports that the two sets of attitudes exist both separately and together, and that the greater the intensity of anti-Israel attitudes, the more likely they are to be accompanied by antisemitic attitudes \citep{staetsky2017}. The elastic view is close in spirit to a latent-trait model, in which a person's position on a continuum, rather than membership of a category, is the primary quantity. Keeping the two attitude domains separate and then relating them is the correlated-factors approach of \cref{sec:model}.

\paragraph{Other general-population surveys}
The American Jewish Committee (AJC) runs parallel annual surveys of American Jews and of the general public. For its 2024 report, both were conducted by SSRS. The general-public survey used the SSRS Opinion Panel, a probability-based panel \citep{ajc2024general,ajc2024jews}. In academic research, \citet{allington2022} proposed the Generalised Antisemitism (GeAs) scale. It is designed to measure antisemitism as expressed in relation both to Jews and to Israel, and its authors describe it as consistent with the IHRA working definition. \Citet{hersh2023} surveyed about 3,500 US adults with two tests: one of overt antisemitic attitudes and one of double standards applied to Jews. They report that, on their measures, overt antisemitic attitudes were more common on the right than on the left, especially among young adults on the far right. On the double-standards test, they report a double standard against Jews on the left but not on the right, where respondents nevertheless attributed more responsibility to Jews for Israel's conduct than respondents on the left did \citep{hersh2023}. Whether items about Israel measure antisemitism is itself disputed (\cref{sec:bg-definitions}), and that dispute is exactly what a test of the measurement model can inform.

\paragraph{Surveys of Jewish people}
A different family of surveys measures the experiences and perceptions of Jewish people. The EU Agency for Fundamental Rights (FRA) has surveyed self-identified Jewish people directly: its second survey (2018) covered 12 EU Member States, and its third, published in July 2024, covered 13 \citep{fra2018antisemitism,fra2024antisemitism}. The Pew Research Center's \emph{Jewish Americans in 2020} surveyed 4,718 Jewish adults between November 2019 and June 2020 and included questions on perceived antisemitism \citep{pew2021jewish}. These surveys estimate how antisemitism is perceived and experienced by its targets. That is a different estimand from the distribution of attitudes in the general population, and we do not use them as measures of the latter. They matter for this paper in two ways. They inform which Israel-related statements Jewish respondents regard as hostile, which is one input to item selection (\cref{sec:design}). And they are a reminder that respondents exposed to our items include Jewish people (\cref{sec:ethics}).

\subsection{Prior evidence on anti-Israel and antisemitic attitudes}\label{sec:bg-prior}

Several studies have related measures of anti-Israel and antisemitic attitudes.
\begin{itemize}
  \item \citet{kaplan2006} analysed a survey of 500 respondents in each of ten European countries. After controlling for many potentially confounding factors, they found that anti-Israel sentiment consistently predicted the probability that a respondent held antisemitic views. They noted that their analysis makes no assumption about, and does not establish, the direction of causation.
  \item \citet{cohen2009} proposed that antisemitism is often expressed as hostility toward Israel. They tested this with experiments that manipulated mortality salience. In their first study, reminders of mortality increased antisemitism and reduced support for Israel, but only when participants believed that their true attitudes could be detected (a ``bogus pipeline'' condition). The authors interpret this as evidence that social desirability can mask such attitudes. In later studies, mortality salience made Israel, and no other country, appear larger, and it increased punitiveness toward human-rights violations attributed to Israel more than toward identical violations attributed to Russia or India.
  \item \citet{beattie2017} studied US participants. He examined how knowledge of conflicts around the world, antisemitism and Islamophobia relate to criticism of the Israeli and Saudi governments. He reports that some criticisms of the Israeli government are more reliable indicators of antisemitic sentiment than others.
  \item The JPR study reports that the two sets of attitudes also occur separately, and a gradient: more intense anti-Israel attitudes are more likely to be accompanied by antisemitic attitudes \citep{staetsky2017}. ADL describes trope belief and anti-Israel sentiment as distinct but overlapping \citep{adl2023attitudes}.
\end{itemize}
Taken together, these studies report positive associations between the two attitude domains. Two of them report that the strength of the association depends on the content of the Israel-related measure \citep{beattie2017,staetsky2017}. The studies differ in population, items, design and model, so their estimates are not directly comparable, and we do not pool them. Three features limit what they can say about the boundary question. First, a positive correlation between two scales is compatible both with a common latent cause and with distinct but correlated attitudes; the item-level decomposition of \cref{def:boundary} is needed to tell these apart, and \cref{thm:ident} states when it can. Second, comparisons across ideological, age or national groups presuppose that items function equivalently in those groups (\cref{sec:bg-invariance}). Third, direct questions on both topics may be subject to social-desirability bias, and \citet{cohen2009} offer experimental evidence that it can matter (\cref{sec:bg-list}).

\subsection{Why measurement invariance matters}\label{sec:bg-invariance}

A comparison of attitude scores between groups is a comparison of the latent attitude only if the items relate to it in the same way in every group. Formally, measurement invariance holds when the conditional distribution of the item responses given the latent variable does not depend on group membership \citep{meredith1993,millsap2011}. In factor-analytic terms there is a hierarchy of increasingly strong conditions. Configural invariance requires the same pattern of loadings. Metric (weak) invariance requires equal loadings. Scalar (strong) invariance adds equal intercepts or thresholds. Strict invariance adds equal residual variances \citep{meredith1993,steenkamp1998,davidov2014}. Comparing latent means requires at least scalar invariance for the items that carry the comparison. Comparing the boundary estimands, which are functions of loadings, requires invariant loadings on both the general and the specific factor.

Item-level violations are called differential item functioning (DIF). Uniform DIF shifts an item's thresholds; non-uniform DIF changes its discrimination. Standard tests include the Mantel--Haenszel procedure, which compares groups matched on total score \citep{holland1988}, and likelihood-ratio tests of item parameters within an item response model \citep{thissen1993}. When some items are non-invariant, group comparisons can still be identified through a subset of invariant \emph{anchor} items. This is partial invariance \citep{byrne1989}. Alternatively, the alignment method estimates group-specific parameters while seeking a solution with as few large non-invariances as possible \citep{asparouhov2014}. Fit-index changes are often used to judge invariance in practice, but their sensitivity depends on which parameters are non-invariant and on group sizes \citep{chen2007}. For cross-national comparisons, where translation and response styles add further sources of non-equivalence, anchoring vignettes offer a design-based correction \citep{king2004}.

Israel-related items are natural candidates for DIF. A statement about a contested political dispute can attract agreement for reasons tied to a group's political identity, religious commitments or generational experience, and those reasons need not run through the general factor. That is a hypothesis to be tested, not a finding. If such DIF exists and is ignored, it biases group comparisons of latent means and of the boundary estimands; \cref{thm:dif-bias} gives the first-order bias. Group comparisons then require anchors, and \cref{prop:anchor} shows that the anchors themselves must be invariant. The same issue affects stereotype-count indices: \cref{prop:index} shows that a difference in the share of a group agreeing with $k$ of $N$ statements need not reflect a difference in the underlying attitude when invariance fails.

\subsection{Why list experiments}\label{sec:bg-list}

Respondents may answer sensitive questions in the way they believe is socially acceptable rather than truthfully. Misreporting on sensitive questions is well documented, and self-administration tends to reduce it \citep{tourangeau2007}. In a randomised comparison of modes, web administration elicited more reporting of sensitive information than conventional telephone interviewing, and whether respondents found an item sensitive depended on their own actual status \citep{kreuter2008}. For questions about Jews and Israel the direction of the bias need not be the same for everyone. Under the social reference theory of \citet{blair2020}, sensitivity bias depends on whom respondents imagine they are answering to and what they believe that reference group expects. Different groups may therefore under-report different responses, which is itself a form of group-dependent measurement error.

The list experiment, or item count technique \citep{miller1984}, protects individual answers by design. Respondents are randomly assigned to a control group, which sees a list of $J$ non-sensitive items, or a treatment group, which sees the same list plus the sensitive item. Each respondent reports only \emph{how many} items apply to them, not which. Under the assumptions of no design effect (adding the sensitive item does not change answers to the control items) and no liars (respondents in the treatment group answer the sensitive item truthfully within the count), the difference in mean counts between the two groups estimates the prevalence of the sensitive response \citep{glynn2013,blair2012}. Multivariate regression estimators allow the prevalence to depend on covariates \citep{imai2011,corstange2009,blair2012}. \Citet{blair2012} propose a statistical test for design effects, and models that adjust for floor and ceiling effects: strategic misreporting that arises when a truthful count at an extreme of the scale would reveal, or appear to reveal, the answer to the sensitive item. Direct and list responses can also be combined. Under a monotonicity assumption that rules out false admissions, the combined estimator is more efficient than the list experiment alone \citep{aronow2015}. Randomised response \citep{warner1965} offers a related form of protection by a different mechanism \citep{blair2015}.

The protection has a cost. List-experiment estimates are much noisier than direct estimates, because the control items add variance; choosing between them is a bias--variance trade-off \citep{blair2020,glynn2013}. \Cref{prop:list-var} quantifies the cost, and \cref{thm:design} chooses $J$ and the allocation of respondents to minimise it. List experiments have been used to study prejudice. \Citet{kuklinski1997} used one to measure white Americans' racial attitudes. They reported racial prejudice to be markedly higher in the South than elsewhere, and argued that social desirability distorts traditional survey measures. \Citet{kuklinski1997ajps} used one to study the relationship between racial prejudice and opposition to affirmative action. They concluded that prejudice is a real problem but cannot by itself explain whites' anger over affirmative action. \Citet{imai2011} reanalysed list-experiment data on racial attitudes from the 1991 National Race and Politics Survey. Evidence on the size of the bias counsels caution. A meta-analysis of list experiments found that sensitivity biases are typically smaller than ten percentage points and are approximately zero in some domains \citep{blair2020}. A correction for social desirability is therefore an empirical question for each item and population, and the design of \cref{sec:design} treats it as one.
